\begin{textsample}{POS Dim 4 – 2015\_c5 – Score 1.00 – 2015\_c5/enade\_2015\_ri\_4d06c54734.txt}  \label{ex:f4_pos_168}
CAPÍTULO 5 DISTRIBUIÇÃO DOS CONCEITOS 5.1 PANORAMA NACIONAL DA DISTRIBUIÇÃO DOS CONCEITOS O Conceito Enade é um indicador de qualidade que avalia o desempenho dos estudantes a partir dos resultados do Enade.
Os valores possíveis vão de 1 ( pior situação ) a 5 ( melhor situação ).
As Unidades de \textbf{Observação} com apenas um ou sem nenhum concluinte participante não obtêm o Conceito Enade, ficando Sem Conceito ( SC ).
A Tabela 5.1 apresenta a quantidade e distribuição de cursos de Relações Internacionais participantes do Enade/2015, por faixa de conceito e Grande Região.
Mais uma vez, ressaltando que a diferença entre os cursos tabulados neste capítulo e no capítulo 2 corresponde aos cursos sem conceito, em princípio, aqueles sem alunos concluintes que participassem da prova ou que tivessem tão somente um aluno concluinte¹⁸.
Observando -se os dados da Tabela 5.1, nota -se que, dos 101 cursos participantes, 39 ( 38,6 \% ) classificaram -se com conceito 3, o valor modal.
Este foi também o conceito modal nas seguintes regiões : Norte ( 60,0 \% ), Nordeste ( 62,5 \% ) e Sudeste ( 41,4 \% ).
Na região Centro-Oeste, os conceitos foram distribuídos de forma bimodal : 33,3 \% dos cursos foram classificados no conceito 3 e 33,3 \% no conceito 2.
Na região Sul, o conceito 4 ( 42,9 \% ) foi o conceito modal.
O conceito 4 foi o segundo mais frequente em nível nacional ( 26,7 \%, correspondendo a 27 cursos ) e o conceito 2, o terceiro ( 20,8 \%, correspondendo a 21 cursos ).
Houve, ainda, oito cursos ( 7,9 \% ) que receberam conceito 5 e quatro ( 4,0 \% ) cursos que receberam conceito 1.
Dos 101 cursos de Relações Internacionais, apenas dois cursos ( 2,0 \% ) ficaram sem conceito ( SC ). ¹⁸ Estes cursos com somente um concluinte não foram considerados no capítulo 2.
A região Norte participou com cinco cursos ou 5,0 \% do total nacional.
Nesta região, o conceito 3 foi atribuído a três cursos dos cinco participantes, o que equivale a 60,0 \% do total regional, além de ser o conceito modal na região.
Nenhum curso ficou sem conceito ( SC ) ou com conceitos 1 ou 5.
Dois cursos foram avaliados com conceitos 2 e 3 ( um curso cada ), correspondendo a 20,0 \%, cada.
A região Nordeste participou com oito cursos ou 7,9 \% do total nacional.
Nesta região, cinco cursos, 62,5 \% em termos regionais, obtiveram conceito 3, o conceito modal para a região.
Os conceitos 4 e 5 foram atribuídos, respectivamente, a um e dois cursos ( 12,5 \% e 25,0 \% ).
Nenhum curso ficou sem conceito ( SC ) ou recebeu conceito 1 ou 2 nessa região.
Dos 58 cursos participantes da região Sudeste ( 57,4 \% do total nacional ), o conceito 3 também foi o conceito modal, e foi obtido por 24 cursos ( 41,4 \% ).
O conceito 1 foi atribuído a dois cursos ( 3,4 \% ) e os conceitos 2 e 4, a 13 e 15 cursos, correspondendo a 22,4 \% e 25,9 \%, respectivamente.
O conceito 5 foi atribuído a três cursos ( 5,2 \% ).
Apenas um curso ficou sem conceito ( SC ).
Os 21 cursos da região Sul ficaram distribuídos em todas as faixas de conceitos.
Há predominância do conceito 4 com 42,9 \% ( conceito modal ), ou seja, nove dos 21 cursos participantes na região Sul.
O conceito 2 e 3 foram atribuídos a oito cursos ( quatro cursos cada ), correspondendo, respectivamente, a 19,0 \% cada um.
O conceito 1 foi atribuído a apenas um curso ( 4,8 \% ) e o conceito 5 foi recebido por dois cursos ( 9,5 \% ).
Somente um curso na região Sul ficou sem conceito ( SC ).
Já dos nove cursos participantes na região Centro-Oeste, três receberam conceito 2 e outros três receberam o conceito 3, sendo, assim, uma distribuição com moda distribuída em dois valores contíguos ( 33,3 \% cada ).
Nesta região, nenhum curso ficou sem conceito ( SC ).
Dos demais cursos, três cursos foram alocados no conceito 3, A região Norte participou com cinco cursos ou 5,0 \% do total nacional.
Nesta região, o conceito 3 foi atribuído a três cursos dos cinco participantes, o que equivale a 60,0 \% do total regional, além de ser o conceito modal na região.
Nenhum curso ficou sem conceito ( SC ) ou com conceitos 1 ou 5.
Dois cursos foram avaliados com conceitos 2 e 3 ( um curso cada ), correspondendo a 20,0 \%, cada.
A região Nordeste participou com oito cursos ou 7,9 \% do total nacional.
Nesta região, cinco cursos, 62,5 \% em termos regionais, obtiveram conceito 3, o conceito modal para a região.
Os conceitos 4 e 5 foram atribuídos, respectivamente, a um e dois cursos ( 12,5 \% e 25,0 \% ).
Nenhum curso ficou sem conceito ( SC ) ou recebeu conceito 1 ou 2 nessa região.
Dos 58 cursos participantes da região Sudeste ( 57,4 \% do total nacional ), o conceito 3 também foi o conceito modal, e foi obtido por 24 cursos ( 41,4 \% ).
O conceito 1 foi atribuído a dois cursos ( 3,4 \% ) e os conceitos 2 e 4, a 13 e 15 cursos, correspondendo a 22,4 \% e 25,9 \%, respectivamente.
O conceito 5 foi atribuído a três cursos ( 5,2 \% ).
Apenas um curso ficou sem conceito ( SC ).
Os 21 cursos da região Sul ficaram distribuídos em todas as faixas de conceitos.
Há predominância do conceito 4 com 42,9 \% ( conceito modal ), ou seja, nove dos 21 cursos participantes na região Sul.
O conceito 2 e 3 foram atribuídos a oito cursos ( quatro cursos cada ), correspondendo, respectivamente, a 19,0 \% cada um.
O conceito 1 foi atribuído a apenas um curso ( 4,8 \% ) e o conceito 5 foi recebido por dois cursos ( 9,5 \% ).
Somente um curso na região Sul ficou sem conceito ( SC ).
Já dos nove cursos participantes na região Centro-Oeste, três receberam conceito 2 e outros três receberam o conceito 3, sendo, assim, uma distribuição com moda distribuída em dois valores contíguos ( 33,3 \% cada ).
Nesta região, nenhum curso ficou sem conceito ( SC ).
Dos demais cursos, três cursos foram alocados no conceito 3, 4 e 5 ( um curso em cada ), correspondendo a 11,1 \% cada um desses três cursos nessa região.
O Gráfico 5.1 apresenta a distribuição do Conceito Enade, segundo Grande Região : poligonais mais à esquerda representam Grandes Regiões com conceituação mais baixa, e poligonais mais à direita, Grandes Regiões com melhores conceitos.
A região Centro-Oeste ( linha preta ) apresenta os cursos com conceitos mais baixos, ao passo que a região Nordeste ( linha verde ) apresenta as distribuições com valores maiores.
Cursos em IES das regiões Norte ( linha azul ), Sudeste ( linha vermelha ) e Sul ( linha roxa ) apresentam uma situação intermediária. 5.2 CONCEITOS POR CATEGORIA ADMINISTRATIVA E POR GRANDE REGIÃO A Tabela 5.2 apresenta a distribuição dos cursos participantes do Enade/2015 de Relações Internacionais, por Categoria Administrativa, de acordo com os conceitos por eles alcançados, segundo as Grandes Regiões brasileiras.
Dos 101 cursos participantes, 22 ( 21,8 \% ) eram ministrados em Instituições Públicas e 79 ( 78,2 \% ), em Instituições Privadas.
De acordo com as informações da Tabela 5.2, em termos nacionais, entre Instituições Públicas e Privadas, nota -se que das oito IES que receberam o conceito 5, seis eram Públicas e duas Privadas.
Dos 22 cursos participantes de IES Públicas, o conceito 4 foi o valor modal, atribuído a nove cursos ( 40,9 \% ).
Entre os demais cursos de Instituições Públicas participantes, apenas um curso ( 4,5 \% ) obteve o conceito 1 e os conceitos 3 e 5 foram atribuídos a seis cursos cada ( 27,3 \% cada ).
Nesta categoria, nenhum curso ficou sem conceito ( SC ).
Na rede privada, o conceito modal foi 3, atribuído a 33 cursos dos 79 da categoria.
Entre os demais cursos participantes, três receberam conceito 1 e 21 cursos obtiveram conceito 2.
O conceito 4 foi atribuído a 18 cursos e o conceito 5, a dois cursos.
Nesta Categoria Administrativa, dois cursos ficaram sem conceito ( SC ).
Na análise por região, observa -se que, na região Norte, as Instituições Privadas participaram com três cursos ( 60,0 \% do total regional ), a menor das proporções desta categoria na região, dos quais nenhum curso ficou sem conceito ( SC ) ou com conceitos 1, 4 ou 5.
O conceito modal para as Instituições Privadas na região foi 3, atribuído a dois Na análise por região, observa -se que, na região Norte, as Instituições Privadas participaram com três cursos ( 60,0 \% do total regional ), a menor das proporções desta categoria na região, dos quais nenhum curso ficou sem conceito ( SC ) ou com conceitos 1, 4 ou 5.
O conceito modal para as Instituições Privadas na região foi 3, atribuído a dois cursos, correspondendo a 66,7 \%, e um curso foi alocado no conceito 2.
As Instituições Públicas participaram com dois cursos na região Norte ( 40,0 \% do total regional ), a maior das proporções desta categoria na região, sendo que os conceitos 3 e 4 foram atribuídos a um curso cada ( 50,0 \%, cada ).
Desse modo, nenhum curso ficou sem conceito ( SC ) ou recebeu os demais conceitos nesta combinação de categoria e região.
Na região Nordeste, a rede privada concentrou cinco dos oito cursos participantes, o equivalente a 62,5 \% do total da região.
Dentre os cursos oferecidos por IES Privadas no Nordeste, quatro cursos ficaram alocados no conceito 3, conceito modal, e um curso obteve conceito 5.
Nenhum curso ficou sem conceito ( SC ) ou recebeu os demais conceitos.
As Instituições Públicas dessa região participaram com três cursos ( 37,5 \% ), verificando -se que os conceitos 3, 4 e 5 foram atribuídos a um curso cada.
Nenhum curso ficou sem conceito ( SC ) ou recebeu os demais conceitos nesta combinação de categoria e região.
Na região Sudeste, a proporção de cursos da rede privada, 84,5 \%, foi mais elevada do que nas demais regiões brasileiras, correspondendo a 49 dos 58 cursos participantes.
Nesta categoria, na região Sudeste, o conceito modal foi 3 ( 21 cursos ).
Nesta combinação de Categoria Administrativa e Grande Região, apenas um curso ficou sem conceito ( SC ) enquanto nenhum recebeu o conceito 5.
Os demais foram avaliados com conceito 1 ( dois cursos ), conceito 2 ( 13 cursos ) e conceito 4 ( 12 cursos ).
A respeito dos nove cursos oferecidos por Instituições Públicas na região Sudeste, sabe -se que a distribuição foi multimodal, visto que os conceitos 3, 4 e 5 foram atribuídos a três cursos cada.
Na região Sudeste, nesta categoria, nenhum curso ficou sem conceito ( SC ) ou recebeu os demais conceitos.
As Instituições Privadas concentraram 15 dos 21 cursos participantes da região Sul, 71,4 \% do total regional.
Desses, cinco ficaram com conceito 4, o conceito modal.
Nesta combinação de Categoria Administrativa e Grande Região, oito cursos ficaram com conceito 2 e 3, quatro cursos cada, enquanto apenas um curso recebeu conceito 5.
Um curso ficou sem conceito ( SC ) e nenhum recebeu o conceito 1.
As Instituições Públicas na região Sul participaram com seis cursos ( 28,6 \% ) e o conceito modal foi 4, atribuído a quatro cursos.
Dois cursos foram avaliados com conceitos 1 e 5 ( um curso cada ).
Nenhum curso ficou sem conceito ( SC ) ou recebeu os conceitos 2 ou 3.
Na região Centro-Oeste, sete dos nove cursos participantes eram de Instituições Privadas ( 77,8 \% em termos regionais ).
Destes, três concentraram -se no conceito 2, conceito modal.
Dentre os demais, dois receberam conceito 3, e dois receberam os conceitos 1 e 4 ( um curso cada ).
Nenhum curso ficou sem conceito ( SC ) ou recebeu o Na região Nordeste, a rede privada concentrou cinco dos oito cursos participantes, o equivalente a 62,5 \% do total da região.
Dentre os cursos oferecidos por IES Privadas no Nordeste, quatro cursos ficaram alocados no conceito 3, conceito modal, e um curso obteve conceito 5.
Nenhum curso ficou sem conceito ( SC ) ou recebeu os demais conceitos.
As Instituições Públicas dessa região participaram com três cursos ( 37,5 \% ), verificando -se que os conceitos 3, 4 e 5 foram atribuídos a um curso cada.
Nenhum curso ficou sem conceito ( SC ) ou recebeu os demais conceitos nesta combinação de categoria e região.
Na região Sudeste, a proporção de cursos da rede privada, 84,5 \%, foi mais elevada do que nas demais regiões brasileiras, correspondendo a 49 dos 58 cursos participantes.
Nesta categoria, na região Sudeste, o conceito modal foi 3 ( 21 cursos ).
Nesta combinação de Categoria Administrativa e Grande Região, apenas um curso ficou sem conceito ( SC ) enquanto nenhum recebeu o conceito 5.
Os demais foram avaliados com conceito 1 ( dois cursos ), conceito 2 ( 13 cursos ) e conceito 4 ( 12 cursos ).
A respeito dos nove cursos oferecidos por Instituições Públicas na região Sudeste, sabe -se que a distribuição foi multimodal, visto que os conceitos 3, 4 e 5 foram atribuídos a três cursos cada.
Na região Sudeste, nesta categoria, nenhum curso ficou sem conceito ( SC ) ou recebeu os demais conceitos.
As Instituições Privadas concentraram 15 dos 21 cursos participantes da região Sul, 71,4 \% do total regional.
Desses, cinco ficaram com conceito 4, o conceito modal.
Nesta combinação de Categoria Administrativa e Grande Região, oito cursos ficaram com conceito 2 e 3, quatro cursos cada, enquanto apenas um curso recebeu conceito 5.
Um curso ficou sem conceito ( SC ) e nenhum recebeu o conceito 1.
As Instituições Públicas na região Sul participaram com seis cursos ( 28,6 \% ) e o conceito modal foi 4, atribuído a quatro cursos.
Dois cursos foram avaliados com conceitos 1 e 5 ( um curso cada ).
Nenhum curso ficou sem conceito ( SC ) ou recebeu os conceitos 2 ou 3.
Na região Centro-Oeste, sete dos nove cursos participantes eram de Instituições Privadas ( 77,8 \% em termos regionais ).
Destes, três concentraram -se no conceito 2, conceito modal.
Dentre os demais, dois receberam conceito 3, e dois receberam os conceitos 1 e 4 ( um curso cada ).
Nenhum curso ficou sem conceito ( SC ) ou recebeu o

% matched lemmas: observação
\end{textsample}
